% Ad M. Brutum 1.14 (ad-brutum-01-14)
% Source: https://raw.githubusercontent.com/PerseusDL/canonical-latinLit/master/data/phi0474/phi059/phi0474.phi059.perseus-lat1.xml
% Fetched by scripts/fetch_latin.py

% Dateline (Perseus): Scr. Romae v Id. Quint. , a. 711 (43).

\ciceroLetterOpener{Cicero Bruto salutem}

\ciceroSection{1} breves litterae tuae—breves dico, immo nullae. tribusne versiculis his temporibus Brutus ad me? nihil scripsisses potius. et requiris meas! quis umquam ad te tuorum sine meis venit? quae autem epistula non pondus habuit? quae si ad te perlatae non sunt, ne domesticas quidem tuas perlatas arbitror. Ciceroni scribis te longiorem daturum epistulam. recte id quidem, sed haec quoque debuit esse plenior. ego autem, cum ad me de Ciceronis abs te discessu scripsisses, statim extrusi tabellarios litterasque ad Ciceronem ut, etiam si in Italiam venisset, ad te rediret; nihil enim mihi iucundius, nihil illi honestius. quamquam aliquotiens ei scripseram sacerdotum comitia mea summa contentione in alterum annum esse reiecta— quod ego cum Ciceronis causa elaboravi tum Domiti, Catonis, Lentuli, Bibulorum; quod ad te etiam scripseram—; sed videlicet, cum illam pusillam epistulam tuam ad me dabas, nondum erat tibi id notum.

\ciceroSection{2} qua re omni studio a te, mi Brute, contendo ut Ciceronem meum ne dimittas tecumque deducas; quod ipsum, si rem publicam cui susceptus es respicis, tibi iam iamque faciendum est. renatum enim bellum est idque non parvum scelere Lepidi. exercitus autem Caesaris, qui erat optimus, non modo nihil prodest sed etiam cogit exercitum tuum flagitari. qui si Italiam attigerit, erit civis nemo quem quidem civem appellari fas sit, qui se non in tua castra conferat. etsi Brutum praeclare cum Planco coniunctum habemus; sed non ignoras quam sint incerti et animi hominum infecti partibus et exitus proeliorum. quin etiam si, ut spero, vicerimus, tamen magnam gubernationem tui consili tuaeque auctoritatis res desiderabit. subveni igitur, per deos, idque quam primum tibique persuade non te Idibus Martiis quibus servitutem a tuis civibus depulisti plus profuisse patriae quam si mature veneris profuturum. v Idus Quintiles.
